\documentclass[10pt]{article}
\usepackage[margin=20mm]{geometry}
\usepackage{amsmath,amssymb,fontspec,unicode-math,xurl,hyperref}
\IfFontExistsTF{DejaVu Serif}{\setmainfont{DejaVu Serif}}{\setmainfont{Latin Modern Roman}}
\setmathfont{latinmodern-math.otf}
\setlength{\parindent}{0pt}\setlength{\parskip}{6pt}\sloppy
\begin{document}
ELVIA Scientific Verification Report
I25-40138 / 0034\par
Paper: Position: AI’s growing due process problem\par
Inspection outcome: Minor: bounded mischaracterization of the Loomis safeguards; other checked claims are supported or need scope clarification\par
Reviewed version: Frozen publisher PDF, PMLR 267:82062-82071 (2025), 10 pages\par
Review date: Review and severity calibration: 24 October 2026 (America/Chicago); machine logs use UTC\par
\section*{1. Identity, verdict and scope}
Sunayana Rane, ICML 2025 position paper [S1]. Page numbers below are the paper’s printed local pages 1-10, which equal one-based PDF pages; they are not the PMLR volume pagination. The frozen original is papers/0034/0034.pdf. Its 156,215 bytes and SHA-256 match the source-acquisition record. No original bytes were changed.\par
One Minor finding is retained after claim-level calibration. The Loomis account omits substantive limits and overstates the absence of information. Correcting it changes the extent of the claimed procedural deficit in that case. The corrected record still leaves the internal scoring weights undisclosed, which is the specific access limitation relevant to the explanation-to-contestation argument; the report does not infer adequacy or inadequacy of the real safeguards from that fact alone.\par
The inspection covered the abstract and Sections 1-7, both principal case studies, the interpretation discussion, alternative views, and selected primary references. It used source comparison and logical analysis, not empirical reproduction. There are no theorem proofs, benchmark tables, or released experimental results in this paper to reproduce. A broad policy position was not relabeled a mathematical impossibility theorem.\par
\section*{2. Definitions and limits that matter}
Four objects must be kept separate: an algorithm’s input-to-score mapping; the reasons an official gives for an adverse decision; the recipient’s practical ability to correct or contest that decision; and the legal adequacy of the entire procedure. A model’s lack of complete mechanistic or causal explanation does not by itself specify all of the other three. Conversely, knowing a model’s coefficients does not by itself establish a fair or legally adequate procedure.\par
The paper itself limits constitutional due process to relevant government action (p. 1) and extends it to private activity as an ideal in Section 6 (p. 8). This distinction is retained. Mathews v. Eldridge, 424 U.S. 319, 334-335 [S4], evaluates procedural sufficiency in context using the private interest, risk of erroneous deprivation and value of extra safeguards, and government interests and burdens. K.W. applies that framework. Neither source supplies a universal requirement that every internal computation be understood. This is a scope clarification, not a claim that a particular opaque deployment is lawful.\par
The audit checks what the cited historical decisions say. It does not adjudicate a live case, establish the current law in every jurisdiction, or certify any AI system. Model code, individual case records, raw COMPAS data, and the underlying Idaho training records were not obtained or executed. No new estimate of bias, recidivism prediction, or benefit adequacy is claimed.\par
\section*{3. Problem 1: the Loomis account loses material conditions}
Source claim and location. Case study 2, p. 6 right column, describes the permission to use COMPAS in sentencing largely through the existence of other reasons for the decision. After reproducing five cautions, p. 7 right column says the court has full discretion whether and how to take the list into account. Section 5, p. 8, summarizes the opposing position as permitting biased systems with a warning label. On p. 7 left column, the accompanying trade-secret discussion says courts have no information about the system whatsoever. These passages together imply a materially weaker set of safeguards and a more complete information vacuum than the primary opinion describes.\par
Decisive check. The official State v. Loomis opinion, 2016 WI 68, filed July 13, 2016 [S2], paragraph 98 says risk scores may not determine incarceration or sentence severity, and may not be the determinative factor in deciding whether an offender can be safely and effectively supervised in the community. Paragraph 99 requires the court to explain independent factors supporting the sentence. Paragraph 100 requires the written advisement and then separately lists the cautions. Thus, independent sentencing reasons and warning labels are not the entire rule. The prohibitions are express conditions of the permission described in paragraphs 8 and 120.\par
The information claim is also too absolute. Paragraphs 51-56 distinguish unavailable factor weights and risk-calculation details from information that is available: the defendant’s report, underlying questions and answers, criminal-history inputs, and portions of the practitioner’s guide. The majority reasons that the defendant can challenge relevant underlying information. Whether this is enough is open to the paper’s criticism, but it is not an absence of all information.\par
Alternative reading and counterevidence. Paragraph 71 gives sentencing courts discretion in assessing a risk tool’s accuracy, and the decision does not prescribe a numerical weight for each caution. The paper can reasonably question whether judges actually avoid cognitive bias or whether the safeguards work. Justice Abrahamson’s concurrence, paragraphs 132-135, expressly records the court’s limited understanding of COMPAS and calls for stronger on-record reasoning; this supports the paper’s concern. These points do not erase the majority’s substantive restrictions or convert partial disclosure into no disclosure. The review therefore does not assert that the court eliminated all discretion, that the warnings are empirically effective, or that the author’s normative disagreement is erroneous.\par
Bounded consequence and verified correction. Replace the blanket account with this statement: Loomis permitted consideration of COMPAS subject to substantive restrictions, independent sentencing reasons and written cautions, while leaving the internal weighting algorithm undisclosed and some input and report information available. The sufficiency and practical effectiveness of those protections can then be criticized explicitly. Each clause is supported by the cited paragraphs. This repairs the factual summary; it does not validate COMPAS, establish absence of bias, or resolve the policy disagreement. The incorrect 2018 date in the paper is not counted as a separate substantive issue.\par
Evidence: source\_pages/source\_p6.png and source\_p7.png visually confirm the paper’s wording; sources/0034/conceptual\_loomis\_2016.pdf preserves the official opinion. The check is primary-source contradiction with a specified alternative reading, not a numerical counterexample or an empirical causal result.\par
Calibration dependency and judgment: Minor retained. The false warning-only/total-information-vacuum account strengthens one case illustration. Its correction still leaves the internal scoring mechanism undisclosed; available inputs permit some challenges without establishing sufficient model-level explanation. Section6, p.8, makes the proposed assessment conditional on technically supported notice and a fair opportunity to contest, and the separately checked K.W. notice/contestation evidence does not depend on Loomis. Thus the source correction alters the claimed extent of one case’s deficit without defeating those independent premises. This bounded factual role, rather than survival of a broad agenda, supports Minor. It does not validate the broadest incompatibility language, establish safeguard efficacy, or certify COMPAS safety or fairness.\par
\section*{4. Checks that did not establish additional defects}
Checked material | Result and boundary\par
K.W. v. Armstrong, pp. 3-4 | The 2016 filed order [S3], pp. 11-15, supports the 733 usable records, major data attrition, underrepresentation, roughly 10-15\% inadequate-budget estimate, and nearly 39\% without a guardian/conservator. These are court-record figures, not independently recomputed population estimates.\par
Secrecy in the Idaho case | The 2016 order, pp. 21-24, specifically handles access to SIB-R assessment material and copyright/confidentiality concerns, including protective access. Earlier trade-secret resistance and disclosure of spreadsheet formulas are described by the plaintiffs’ lawyer in 2023 Senate testimony [S5], p. 3. The testimony supports the historical narrative but is not the earlier court order itself. Its absence from the 2016 order is not evidence that the event never happened.\par
Causal explanation and regression, pp. 2-5 | Coefficients require attention to units, coding, feature values and correlations; they are not automatically causal effects. No numerical coefficient ranking or causal estimator is supplied here that can be falsified. Treat the discussion as advice needing precise implementation, not an invalid derivation.\par
Neural networks, LLMs and trees, pp. 5-6 | The limited point that fluent explanations or saliency displays need not faithfully explain a decision is distinct from a universal claim that every explanation is meaningless. Deeper models do not universally improve out-of-sample accuracy. The paper supplies qualitative cautions rather than a testable scaling law or a proof of universal interpretability failure.\par
Broad incompatibility and alternatives, pp. 3-8 | Section 4 admits different levels of salvageability; p. 4 rejects banning statistical tools as such; Section 6 presents a conditional notice/hearing checklist. These qualifications favor a normative, deployment-specific reading over a literal theorem prohibiting every use of every AI model.\par
Usefulness versus determinative effect, p. 6 | An evidentiary item can corroborate other information or alter confidence without changing a particular final action. The rhetorical either/or should not be read as a formal proof. This clarification does not establish that COMPAS is useful or that judges avoid anchoring.\par
The case study supports serious process failures in the specified system. It does not furnish a controlled comparison of all modern AI systems with human decision-makers. No unmeasured efficacy or prevalence claim has been promoted to a verified empirical result.\par
Abstract-level checks of Adebayo et al. [S6] and Turpin et al. [S7] corroborate limitations of some saliency methods and possible unfaithfulness of generated explanations. They do not prove universal failure. Their experiments were not rerun and these checks do not establish the adequacy of any deployed legal procedure.\par
\section*{5. Dependence, versions and remaining evidence}
As of 24 October 2026, a fresh download of the publisher-linked PDF is byte-for-byte identical to the frozen original (same 156,215 bytes and SHA-256); exact extracted text also matches. The citation gives ICML 13-19 July 2025, while the PMLR repository metadata records 6 October 2025 as its issue/post date. No separate arXiv version or labeled erratum was located in the bounded search. A related earlier CS\&Law poster title was found, but no manuscript was retrieved and it is not treated as a verified prior version. OpenReview TEkyydR6il presented a browser-verification challenge, so its discussions and attachments were not inspected. Problem 1 remains in the current publisher bytes.\par
Problem 1 affects the legal description supporting the paper’s critique of the criminal-justice case. The broader research agenda remains available: improve accessible explanations, distinguish tool types, measure actual harms, permit meaningful challenge, and examine appropriate limits on deployment. The verified correction is a source-faithful rewrite only. Fresh empirical work would be needed to determine whether the Loomis safeguards change judges’ behavior, and deployment-specific legal and technical evidence would be needed to judge a particular system’s procedural adequacy.\par
No later case history was used to retroactively reject the paper’s 2025 account. A complete legal citator search, the entire K.W. docket, all legal and technical references, and contemporary model-specific interpretability capabilities were outside scope. No dedicated code or separate supplement was located in the bounded official-source search; absence of an artifact is not a defect for this position paper.\par
\section*{6. Conclusion and decision ledger}
Decision item | Status and scope\par
Issue judgment | F1: Minor retained. Correct the extent of the Loomis deficit; the corrected record retains the specific undisclosed-mechanism premise, and K.W. evidence is independent.\par
Repair verification | Replacement checked against majority paragraphs 51-56 and 98-100. Safeguard efficacy and COMPAS safety are not established.\par
Author adoption | Not requested. No correction has been represented as accepted.\par
Other inspected claims | Idaho case: No Issue Found within checked scope. Technical and normative claims: scope clarifications, not additional demonstrated defects.\par
Remaining evidence | Earlier Idaho disclosure order and related poster manuscript not retrieved. Judicial-behavior experiments and deployment-specific analysis would be needed for stronger claims.\par
Scientific decision | Pending human adjudication and publication status.\par
\section*{7. Consulted references}
[S1] Rane, S. Position: AI’s growing due process problem. Proceedings of ICML 2025, PMLR 267:82062-82071. Frozen 10-page publisher PDF; accessed 24 October 2026. https://proceedings.mlr.press/v267/rane25b.html\par
State v. Loomis, 2016 WI 68, 881 N.W.2d 749, 371 Wis. 2d 235. Supreme Court of Wisconsin, July 13, 2016. Official 67-page opinion PDF, including concurrences; majority paragraph identifiers control. https://www.wicourts.gov/sc/opinion/DisplayDocument.pdf?content=pdf\&seqNo=171690\par
K.W. v. Armstrong, 180 F. Supp. 3d 703 (D. Idaho 2016). Case 1:12-cv-00022-BLW, Document 270, filed March 28, 2016, 37 pages. Filed judicial opinion preserved in an ACLU Idaho mirror; court-stamped PDF, not an advocacy summary. https://www.acluidaho.org/app/uploads/2016/03/Kw\_v\_armstrong\_decision\_03292016.pdf\par
Mathews v. Eldridge, 424 U.S. 319 (1976), especially 334-335. U.S. Reports scan preserved by the Library of Congress; 32 pages. https://tile.loc.gov/storage-services/service/ll/usrep/usrep424/usrep424319/usrep424319.pdf\par
Eppink, R. Testimony on AI in Government. U.S. Senate Committee on Homeland Security and Governmental Affairs, May 16, 2023, especially p. 3. First-person testimony by counsel, hosted by the Senate; not a judicial finding. https://www.hsgac.senate.gov/wp-content/uploads/Testimony-Eppink-2023-05-16-1.pdf\par
Adebayo, J., Gilmer, J., Muelly, M., Goodfellow, I., Hardt, M., and Kim, B. Sanity Checks for Saliency Maps. NeurIPS 31 (2018). Official abstract inspected only. https://proceedings.neurips.cc/paper/2018/hash/294a8ed24b1ad22ec2e7efea049b8737-Abstract.html\par
Turpin, M., Michael, J., Perez, E., and Bowman, S. R. Language Models Don’t Always Say What They Think: Unfaithful Explanations in Chain-of-Thought Prompting. arXiv:2305.04388v2, 9 December 2023; NeurIPS 2023. Abstract inspected only. https://arxiv.org/abs/2305.04388v2\par
\section*{Supporting checks}
The original, court opinions and Mathews scan were preserved separately.\par
Access limit: Senate testimony p. 3 was read through web retrieval; the local request returned HTTP 403. Its URL, attributed evidence note and failure were saved; no local full Senate PDF is claimed.\par
Public-source re-fetching can obtain different future bytes; saved hashes define this review.\par
No human scientific review or author adoption is claimed. No author code ran. Scientific decisions remain pending.\par
Focused severity calibration 001 (24 October 2026) applies the user’s clarified core-dependency criterion. Existing arithmetic, sources and version checks were reused. Canonical reports, findings, coverage, provenance and checkpoint are synchronized. No new empirical result or author-code execution is claimed.\par
\end{document}
